\documentclass[11pt,a4paper]{article}
\usepackage[T1]{fontenc}
\usepackage[utf8]{inputenc}
\usepackage{lmodern}
\usepackage{amsmath,amssymb}
\usepackage[margin=25mm]{geometry}
\usepackage{longtable,booktabs,array,calc}
\usepackage[hidelinks]{hyperref}
\hypersetup{pdftitle={Corpus audit of 1 October 2026: two-axis scores for all 175 notes, the uncorrected errors },pdfauthor={navstokgap research working notes}}
\newcommand{\tightlist}{\setlength{\itemsep}{0pt}\setlength{\parskip}{0pt}}
\setlength{\emergencystretch}{3em}
\setcounter{secnumdepth}{0}
\begin{document}
\section{Corpus audit of 1 October 2026: two-axis scores for all 175
notes, the uncorrected errors they expose, and where the drift
sits}\label{corpus-audit-of-1-october-2026-two-axis-scores-for-all-175-notes-the-uncorrected-errors-they-expose-and-where-the-drift-sits}

\textbf{Result.} Every maintained note was read in full and scored from
1 to 10 for interestingness and for correctness of its own derivations
by Claude Opus 5.5 under one fixed rubric; Claude Fable 5.1 then
reviewed the condensed statements of the 26 notes scored 6 or more for
interest against their sources and the trap list of \texttt{LLM.md} §4.
Three findings matter for the programme.

\begin{enumerate}
\def\labelenumi{\arabic{enumi}.}
\tightlist
\item
  \textbf{The spine is sound.} The Newton recording results, the
  refinement and gauge insertion results and the strong-coupling
  constants score 8 or 9 for correctness; the reviewer found one false
  claim among the 26 selected statements (item (a) below) and no
  reintroduced trap.
\item
  \textbf{Twenty-four notes carry uncorrected overclaims or inconsistent
  numbers} (correctness 6 or below, four of them at 4). Twenty of the
  twenty-four belong to the Hamiltonian, flow and large-field mass-gap
  attack that preceded the 17 September change of direction, routes that
  \href{mass-gap-position.md}{the position note} §3 lists as closed. The
  closure is recorded in the map; the route notes' own texts keep the
  claims that the closure withdrew. That is the drift: supersession
  lives in the map and in \texttt{LLM.md}, while the superseded notes
  still read as live results.
\item
  \textbf{Interest is concentrated.} The scorer put 26 notes at 6 or
  above and none above 7; 70 of 175 notes are maps or calculations (kind
  counts below). The distributions are a calibration of what an outside
  expert would take from the corpus, under the rubric stated in §1, and
  settle nothing about novelty.
\end{enumerate}

The two digests built from the scores, \emph{The most interesting
results} (26 notes) and \emph{Interesting and correct results} (the 22
of them scored 8 or more for correctness), are published on the
director's personal site as phone-format PDFs with a landing page that
tabulates every score:
\url{https://lxbifi11.bifi.unizar.es/navstokgap/}. The repository keeps
the scores here, in §4, and the corrections to make, in §3.

\subsection{1. Method and reading
labels}\label{method-and-reading-labels}

Six Opus readers worked in parallel on alphabetical batches of about
62,000 words each, each note read line by line, with no scripts for
mathematics (project rule). The rubric's anchors: interest 9 or 10 for a
theorem with a surprising conclusion or a precise closure of a natural
route, 7 or 8 for a solid nontrivial theorem, construction or decisive
obstruction with explicit constants, 5 or 6 for useful supporting work,
below that for bookkeeping; correctness 9 or 10 for complete proofs at
the stated level with every hypothesis stated where used, 7 or 8 for
essentially correct with minor gaps, 5 or 6 where a step is unjustified
or claims outrun proofs, 4 or below for a substantive error or an
unmarked superseded claim. Correctness was scored on the current text,
so a note that was wrong and now carries a dated correction is scored on
the corrected text. Each reader also wrote a 60 to 260 word condensed
statement per note and a one-line reason for each score; the reasons are
quoted in §3 for the low scores.

The reviewer (Fable) read the 26 selected statements against the notes'
lead sections and theorem statements, checked the arithmetic of the
constants that appear in them, and checked them against the 26 traps of
\texttt{LLM.md} §4. Items marked \emph{verified} below were checked by
the reviewer from the note's own formulas; items marked \emph{reported}
are the scorer's reasons, quoted and unverified here. Reading label for
this note: scorer full-read of every note; reviewer passage-level reads
of the 26 selected notes and grep-level checks of the notes in §3.

Scores are a model's reading with reasons on file. They are not
refereeing, and a high correctness score means a complete derivation at
the level the note states, which may be conditional, formal or a
reading.

\subsection{2. Distributions}\label{distributions}

Interest: 1: 3, 2: 8, 3: 41, 4: 58, 5: 39, 6: 21, 7: 5. Correctness: 4:
4, 5: 6, 6: 14, 7: 27, 8: 59, 9: 65. Kind: calculation: 44,
construction: 13, exploratory: 5, map: 26, obstruction: 24, reading: 12,
theorem: 51. Status as read by the scorer: conditional: 18, conjectural:
4, mixed: 39, none: 16, proved: 96, superseded: 2.

Selected for the digests (interest 6 or more):
\texttt{mark-cost-and-statistical-floor}, \texttt{planck-gap-paper} and
\texttt{planck-gap-probabilistic} at 7; the other 23 at 6. Of the 26,
the four below 8 for correctness are \texttt{planck-gap-paper} (7),
\texttt{four-dimensional-composition} (7),
\texttt{su2-midplane-small-field} (7) and \texttt{planck-gap-derivation}
(6).

\subsection{3. Errors and overclaims
found}\label{errors-and-overclaims-found}

\subsubsection{3a. Verified by the
reviewer}\label{a.-verified-by-the-reviewer}

\begin{enumerate}
\def\labelenumi{(\alph{enumi})}
\item
  \textbf{\texttt{planck-gap-derivation}, boxed area forms off by a
  factor 2.} The note's theorem \(F^2\tau^3>9m\kappa\) with
  \(\tau\Delta E=F^2\tau^3/(2m)=(3F/v)A\) and \(A=vF\tau^3/(6m)\) gives
  \(A>3v\kappa/(2F)\); the box states \(A>3v\kappa/F\). Corollary 3's
  two-mark upper value is likewise \(6v\kappa/F\), stated as
  \(12v\kappa/F\). The theorems themselves (constants 9 and 36) hold.
  The digest on the site carries the corrected values and a reviewer's
  note. The interval model is in any case superseded by the Gaussian
  mark model of
  \href{mark-cost-and-statistical-floor.md}{mark-cost-and-statistical-floor}
  and \href{planck-gap-paper.md}{the Planck paper}, where
  \(\kappa\ge\hbar/2\) holds for every probe state, and the note should
  open by saying so.
\item
  \textbf{\texttt{strong-coupling-target-box}, the \(q=4\) tolerance
  column disagrees with the note's own criterion.} From
  \(e\,\eta\,2^q\le0.15g^2-384e/g^2\) (the note's own criterion) one
  gets \(\eta\le0.105\) at \(g^2=100\) and \(\eta\le0.570\) at
  \(g^2=200\); the table prints \(0.29\) and \(0.62\). The \(q=8\)
  column checks.
\item
  \textbf{\texttt{small-field-step-decay-and-threshold}, abstract and
  body disagree on constants.} The abstract quotes an error density
  \(10^{-12}\); §3 derives \(1.1\times10^4\) and \(10^{-13}\). The
  scorer also reports that §3 sums an \(\ell^\infty\) decay with the
  \(\ell^1\) product formula, undercounting by about 24; that step is
  reported, unverified here.
\end{enumerate}

\subsubsection{3b. Reported by the scorer, with the scorer's
reason}\label{b.-reported-by-the-scorer-with-the-scorers-reason}

\begin{itemize}
\tightlist
\item
  \texttt{flow-instability-large-field} (I 5, C 4): Proposition 1
  (symmetry of \(M\)) is correct. The sharpness claim fails on the
  note's own formula: with the Landau level \(gB\) included,
  \(\omega^2=k_\parallel^2-gB\), so the growth rate is \(gB\), half the
  Duhamel exponent \(2\|G\|\). A symmetric operator of norm \(\|M\|\)
  need not attain \(+\|M\|\). Section 4's `sharp' crossover inherits
  this; the text carries no correction.
\item
  \texttt{large-field-action-lower-bound} (I 3, C 4): Proposition 1 is
  correct in the continuum. Proposition 2 fails: a field constant on the
  block and zero outside violates the Bianchi identity, and flowing at
  radius equal to the block size changes the block integral by order
  one, not \(O(a/\ell)\). Section 3's descent argument is a non
  sequitur. Section 4 imports the unsupported `sharp \(e^{2\eta}\)'
  claim.
\item
  \texttt{mass-gap-conditional-theorem} (I 4, C 4): The decay-to-gap
  transfer and rate bookkeeping in §3 are fine. But §4's `nothing else
  is missing' and `a proof consists of H1 and H2' overclaim: continuum
  existence, OS axioms and nontriviality are not implied. Existence of
  the renormalized interaction and gauge-theory strong mixing are
  glossed. The later rejection of physical-coupling H2 certification is
  not marked.
\item
  \texttt{schur-error-ultraviolet} (I 5, C 4): Lemma 1' and the
  oscillator derivative check. Section 4's own sum gives theta
  \textasciitilde{} g\^{}2 (Lambda L)\^{}3, an extensive vacuum-shift
  slope, not the stated g\^{}2 (Lambda L)\^{}2, so the regime g
  \textless\textless{} 1/N\_s in Section 5 appears wrong (g
  \textless\textless{} N\_s\^{}\{-3/2\}). The Berry coefficient seems to
  drop the four charge/polarization copies. The fixed-lattice theorem is
  asserted but unproved.
\item
  \texttt{gapped-set-critical-coupling} (I 4, C 5): Proposition 1 is
  fine. Proposition 2's proof is heuristic: it equates finite-volume gap
  closing with infinite-volume degeneracy or diverging length, and
  ignores massless phases that fill intervals of coupling. Section 4
  calls the compact \(U(1)\) transition second order with a Maxwell
  limit at \(g_c\), against lattice evidence of weak first order. The
  dated Wilson \(SU(N\ge5)\) precision is present.
\item
  \texttt{ground-state-measure-transfer} (I 4, C 5): Proposition 1 is
  standard and correct. The claim that isotropic Wilson-measure
  estimates hold `verbatim' ignores that the Hamiltonian limit is the
  anisotropic \(a_t\to0\) measure. The bound
  \(\mathbb P\lesssim e^{-c\eta^2/g(\ell)^2}\) in section 3 is a
  heuristic presented as available. The section 4 competition is marked
  superseded in part; the action display carries a stray factor.
\item
  \texttt{large-field-entropy-count} (I 4, C 5): The Landau degeneracy,
  the longitudinal disc count, \(\mathcal N=\eta^2/(8\pi^2)\) and the
  ratio \(2\pi^2/g^2\) are arithmetically correct. The convergence
  claims in §§3--4 (`settles', sum converges for \(g^2<2\pi^2\)) rest on
  an unproved bounded-factor-per-mode premise and on counting only
  constant-background modes. Proposition 1 is trivial.
\item
  \texttt{mass-gap-position} (I 4, C 5): Internal numbers disagree: the
  strong threshold appears as g\^{}2\textgreater=176, \textgreater444
  and \textgreater=1056; the weak side as
  1/g\textsuperscript{2\textgreater\textasciitilde10}2 and as
  g\textsuperscript{2\textasciitilde10}-12. `Large-field region closed
  in measure' exceeds the cost and entropy ingredients cited. `Exhaust
  the methods built from operator norms' generalizes from five specific
  closures.
\item
  \texttt{strong-coupling-target-box} (I 5, C 5): The criterion e eta
  2\^{}q \textless= 0.15 g\^{}2 - 384e/g\^{}2 and the q=8 column check.
  The q=4 column does not: at g\^{}2=100 the criterion gives 0.105, the
  table 0.29; at 200 it gives 0.57, the table 0.62. Tolerances `divided
  by 24 at the listed couplings' fail, since 24*384e/g\^{}2 exceeds
  0.15g\^{}2 there. The Proposition claims more than the cyclic-subspace
  criterion.
\item
  \texttt{typical-field-strength-window} (I 4, C 5): The free-field
  integral, the range condition \(\kappa^2\ge g/4\pi+1/2\) and the
  Gaussian tail estimate check, with \(\hbar\) absorbed into \(g^2\).
  The claims that the Hamiltonian route terminates and that no measure
  exists ignore the vacuum state's configuration distribution. A top box
  flags the lattice correction, but the title, section 4 and section 5
  still state the superseded conclusion.
\item
  \texttt{abelian-misses-the-box} (I 4, C 6): b\_0=11/(16 pi\^{}2) and 2
  b\_0 log 2=0.0966 check. The U(1) rerun of the expansion is asserted
  `verbatim' without the threshold computed. `Gapped set is
  (g\_c,infinity)' (intro, §2) overclaims: no monotonicity or single
  transition is proved. `All three isolations are b\_0 statements' (§3)
  is an unlabelled reading.
\item
  \texttt{action-floor-yang-mills-gap} (I 5, C 6): Prop. 1, Thm 2, Prop.
  4 and the Bohr-Sommerfeld levels check. Remark 3 counts only the
  x-valley (both valleys double the constant; growth law unaffected).
  Prop. 5's d=4 row drops 1/g\^{}2, which by the table's own convention
  has action dimension, so `no constant at all' in d=4 is inconsistent.
\item
  \texttt{agmon-ground-state-suppression} (I 3, C 6): Lemma 1 and
  Corollary 2 are correct; Corollary 2 with rho=(1-delta)d needs epsilon
  tied to delta, left implicit. The distance step (§2) is a labelled
  scaling argument. The correction box flags that only the global excess
  is controlled, yet §3-§4 still state local consequences (finite
  exponential local moments) that do not follow.
\item
  \texttt{confinement-scale-bands} (I 5, C 6): Table entries (r\_0/a, m
  a=4.21/(r\_0/a), sigma a\textsuperscript{2=1.35/(r\_0/a)}2),
  g\^{}2=6/beta, the 444 threshold and the area-law values check; a
  dated caveat flags the softened glueball near beta=5.7. Treating a
  coarse Wilson loop as polymer activity (§2b) is heuristic. `Six
  blocking steps is the whole problem' overstates: the weak side's
  `perturbative in kind' doublings are unproved.
\item
  \texttt{flow-before-decimation} (I 4, C 6): The invariance proposition
  is correct as stated. The sweeping claim that every norm-based
  criterion is exhausted outruns it, since a conjugated Hamiltonian can
  be split differently. The addendum's local \(e^{-\lambda n}\)
  ground-state suppression is superseded by the later Agmon-global
  correction and is not marked as such here.
\item
  \texttt{flow-jacobian-truncation-error} (I 5, C 6): Proposition 2's
  Kato-Duhamel bound is correct. Proposition 1 handles the gauge-fixing
  term loosely (`to the same order'). The truncation error is a
  kernel-tail heuristic, not an operator relative bound. The Sharpness
  box repeats the instability note's unsupported claim of growth rate
  exactly \(2\|G\|\). The lattice-step pessimism is correctly flagged as
  corrected.
\item
  \texttt{intermediate-region-finite-verification} (I 5, C 6): The
  decay-to-gap step via reflection positivity is sound. The lead's claim
  that the problem `reduces to' two numbers is weakened only in §4b. The
  lead says gauge invariance is no obstruction, while §5 admits Elitzur
  empties the single-link form; the tension is left unresolved. Box
  sizes are corrected in a later note, flagged here. Literature is at
  metadata level.
\item
  \texttt{large-field-operator-inequality} (I 5, C 6): Propositions 1
  and 2 (min-max and the ground-state-transform trial state) are
  correct, and the IMS arithmetic giving \(C_L=32\pi^2c_0^2/3\) checks.
  The link-derivative bound behind \(|\nabla f_B|^2\) is heuristic. The
  claim that every Hamiltonian route loses a surface term is shown for
  one crude trial construction only; window constants are crude.
\item
  \texttt{lattice-truncation-uniform} (I 5, C 6): Coupling cancellation
  in the lattice flow and the compactness bound are correct, and
  exponential locality of the linearized flow over unit time is
  standard. The Proposition's proof is a sketch: a kernel tail is
  equated with a relative operator error, the Jacobian potential from
  \(\rho_t\) is not treated, and `K=3 suffices' is asserted without
  values for \(C,c\).
\item
  \texttt{planck-gap-derivation} (I 6, C 6): Lemmas 1-3, Theorem 1
  (constant 9) and Theorem 2 (36) check. The area forms are off by a
  factor 2: since (3F/v)A = tau Delta E, the bound is A \textgreater{}
  3v kappa/(2F), and Corollary 3's upper value is 6v kappa/F. The note
  omits the later reinstatement of a quantum kappa = hbar/2.
\item
  \texttt{small-field-step-decay-and-threshold} (I 5, C 6): The
  soft-constraint gap 2/9 at a\textgreater=130 and the rate 1.7e-3
  check, as do the strip facts (a),(b) of §2 for a one-direction shift.
  §3 sums an l-infinity decay with the l1 product coth\^{}4,
  undercounting by about 24. The abstract quotes 2e4 and 1e-12, while §3
  and its table give 1.1e4 and 1e-13. `Closed' outruns a heuristic
  estimate.
\item
  \texttt{strong-coupling-threshold-explicit} (I 5, C 6): The
  Kotecky--Preiss chain checks: degree 80, c\textasciitilde220,
  rho\textless=8e-4, t\_0\textgreater=7.1, e\^{}\{64 t\_0\},
  beta\_*\textasciitilde e\^{}\{-465\}, g\_0\^{}2\textasciitilde1e101.
  §3's 40--70 is an estimate that multiplies the per-site
  beta=54/g\^{}4, which already sums three plaquettes, by a further
  coordination 12. The text does not point to the later rigorous 388/79.
  The weak-side figure 1/g\^{}2\textasciitilde1e2 predates the
  small-field note's 1e-12.
\item
  \texttt{su3-constants} (I 3, C 6): b\_0=11/(16 pi\^{}2), 51/121,
  54/g\^{}4, the 8/3 g\^{}2 flux loop, v\textless=96e c/g\^{}2, Haar
  variance 1/2, the Delta V identity and the doubling table check. The
  blocking paragraph writes beta\_direct=24/g\^{}4; the cited note's
  32N/(C\_2 g\^{}4) gives 72/g\^{}4, which is needed for 3M\^{}2/8.
  `Equivalence of T2' with no zero-temperature transition' is asserted
  as standing without proof.
\item
  \texttt{what-would-unblock} (I 4, C 6): Most citations and
  characterizations are accurate at metadata level. The infrared-bound
  paragraph gives the wrong reason: Froehlich--Simon--Spencer cover O(N)
  spins, so failure stems from the non-quadratic plaquette interaction.
  Section 6 states as fact that coarse effective interactions enter the
  completely analytical set. Item 1 is half-retracted in place while the
  opening still says `exactly two things'.
\end{itemize}

\subsubsection{3c. Where the drift sits}\label{c.-where-the-drift-sits}

Of the 24 notes above, 20 are mass-gap route notes written before 17
September (Hamiltonian flow, large-field region, Agmon, blocking,
truncation, Schur and threshold notes);
\href{mass-gap-position.md}{mass-gap-position} §3 closes their routes
and \texttt{LLM.md} §4 records the traps, but the notes themselves still
state their conclusions in the lead and keep inconsistent thresholds
(176, 444 and 1056 for the Wilson strong-coupling gap appear with
different labels across \texttt{mass-gap-position},
\texttt{su3-constants} and \texttt{strong-coupling-threshold-explicit}).
The other four are \texttt{planck-gap-derivation} (item (a)),
\texttt{action-floor-yang-mills-gap} (a \(d=4\) row that drops
\(1/g^2\)), \texttt{abelian-misses-the-box} (a gapped set stated without
monotonicity) and \texttt{what-would-unblock} (a wrong reason for the
failure of infrared bounds). The repository's own convention for a
withdrawn claim is a dated correction box at the top of the note;
applying it to the 24 notes, with the unified thresholds, would remove
the drift without deleting the record.

\subsection{4. Scores for every note}\label{scores-for-every-note}

Interest (I) and correctness (C) from 1 to 10; status and kind as read
by the scorer.

\begin{longtable}[]{@{}
  >{\raggedright\arraybackslash}p{(\linewidth - 8\tabcolsep) * \real{0.2000}}
  >{\raggedleft\arraybackslash}p{(\linewidth - 8\tabcolsep) * \real{0.2000}}
  >{\raggedleft\arraybackslash}p{(\linewidth - 8\tabcolsep) * \real{0.2000}}
  >{\raggedright\arraybackslash}p{(\linewidth - 8\tabcolsep) * \real{0.2000}}
  >{\raggedright\arraybackslash}p{(\linewidth - 8\tabcolsep) * \real{0.2000}}@{}}
\toprule\noalign{}
\begin{minipage}[b]{\linewidth}\raggedright
Note
\end{minipage} & \begin{minipage}[b]{\linewidth}\raggedleft
I
\end{minipage} & \begin{minipage}[b]{\linewidth}\raggedleft
C
\end{minipage} & \begin{minipage}[b]{\linewidth}\raggedright
Status
\end{minipage} & \begin{minipage}[b]{\linewidth}\raggedright
Kind
\end{minipage} \\
\midrule\noalign{}
\endhead
\bottomrule\noalign{}
\endlastfoot
abelian-misses-the-box & 4 & 6 & mixed & map \\
action-floor-yang-mills-gap & 5 & 6 & mixed & theorem \\
action-scale-dilation & 4 & 9 & proved & obstruction \\
action-scale-obstructions & 4 & 8 & mixed & map \\
action-unit-dimensional-selection & 5 & 9 & proved & theorem \\
additive-noise-marks & 5 & 9 & proved & theorem \\
agmon-global-not-local & 5 & 8 & proved & obstruction \\
agmon-ground-state-suppression & 3 & 6 & mixed & calculation \\
ancient-cuts-provenance & 2 & 8 & none & reading \\
arrow-not-sling & 4 & 7 & none & reading \\
autonomous-finite-readout & 5 & 9 & proved & construction \\
block-apparatus-composition & 3 & 8 & proved & calculation \\
blocking-criterion-monotone & 5 & 7 & proved & obstruction \\
blocking-step-obstruction & 3 & 7 & superseded & exploratory \\
bound-orbit-action-observable & 3 & 9 & proved & calculation \\
bounded-acceleration-return & 3 & 9 & proved & calculation \\
bridge-crossover & 1 & 8 & none & map \\
calibrated-canonical-ambiguity & 3 & 8 & proved & construction \\
calibrated-displacement-ambiguity & 4 & 9 & proved & obstruction \\
calibration-tolerance-recovery & 4 & 9 & proved & theorem \\
causal-force-information & 3 & 9 & proved & calculation \\
checkerboard-dynamics & 4 & 9 & proved & calculation \\
classical-cut-state & 4 & 9 & proved & calculation \\
classical-orientation-closure & 4 & 9 & proved & obstruction \\
classical-readout-refinement & 3 & 9 & proved & construction \\
clock-position-local-recovery & 3 & 9 & proved & theorem \\
closed-orbit-force-action & 4 & 9 & proved & theorem \\
comparison-and-bridges & 3 & 9 & none & map \\
composition-crossover-gap-checks & 3 & 7 & mixed & calculation \\
composition-universality & 4 & 9 & conditional & theorem \\
cone-time-refinement & 2 & 8 & none & map \\
confinement-scale-bands & 5 & 6 & mixed & map \\
conservative-harmonic-receiver & 4 & 8 & mixed & calculation \\
correlated-calibration-response & 3 & 8 & conditional & calculation \\
cut-measure-newton & 6 & 9 & proved & theorem \\
cut-paradox-two-faces & 5 & 8 & proved & theorem \\
cut-point-consistency & 4 & 9 & proved & calculation \\
dimension-ladder & 5 & 7 & conjectural & map \\
discrete-substrate-actors & 3 & 7 & none & exploratory \\
dobrushin-uniqueness-wilson & 5 & 9 & proved & theorem \\
energy-constrained-apparatus-ambiguity & 3 & 9 & proved & obstruction \\
energy-depot-action-selection & 4 & 9 & proved & calculation \\
final-clock-momentum-recovery & 3 & 9 & proved & theorem \\
finite-depth-spin-gap & 2 & 9 & proved & calculation \\
finite-horizon-minimax & 4 & 9 & proved & theorem \\
finite-precision-cut & 3 & 9 & proved & calculation \\
finiteness-half-flowed-susceptibility & 5 & 7 & conditional & theorem \\
fixed-coupling-calibration & 4 & 8 & proved & calculation \\
fixed-force-small-circles & 5 & 9 & proved & obstruction \\
fixed-preparation-ambiguity & 5 & 9 & proved & theorem \\
flow-before-decimation & 4 & 6 & mixed & obstruction \\
flow-conjugation-truncation & 5 & 7 & mixed & obstruction \\
flow-instability-large-field & 5 & 4 & mixed & obstruction \\
flow-jacobian-truncation-error & 5 & 6 & mixed & calculation \\
flowed-bound-free-field & 4 & 9 & proved & calculation \\
four-dimensional-composition & 6 & 7 & mixed & theorem \\
four-dimensional-parallel-log & 6 & 8 & conditional & calculation \\
full-apparatus-preparation-ambiguity & 5 & 9 & proved & theorem \\
full-clock-phase-recovery & 4 & 8 & proved & calculation \\
full-pointer-recovery & 4 & 8 & proved & calculation \\
galileo-two-path-interference & 6 & 9 & conditional & theorem \\
gapped-set-critical-coupling & 4 & 5 & mixed & map \\
gaussian-blocking-coupling & 4 & 8 & proved & calculation \\
global-clock-speed-ambiguity & 4 & 8 & proved & construction \\
ground-state-measure-transfer & 4 & 5 & mixed & exploratory \\
halving-atlas & 4 & 8 & mixed & map \\
hamiltonian-finite-closure & 4 & 9 & proved & obstruction \\
hamiltonian-moment-descent & 4 & 9 & proved & construction \\
hidden-clock-ambiguity & 4 & 9 & proved & construction \\
holography-lowest-dimensions & 3 & 7 & none & reading \\
i003-double-limit-rigidity & 3 & 7 & conjectural & exploratory \\
indistinguishable-phase-bound & 4 & 9 & proved & theorem \\
interacting-ising-gap & 3 & 9 & proved & theorem \\
intermediate-region-finite-verification & 5 & 6 & conjectural & map \\
ising-foundational-value & 2 & 8 & none & map \\
ising-hermitian-transfer & 3 & 9 & proved & construction \\
jacobi-kernels-distinguishability & 1 & 8 & none & map \\
kogut-susskind-strong-coupling-explicit & 6 & 8 & proved & theorem \\
large-field-action-lower-bound & 3 & 4 & mixed & calculation \\
large-field-entropy-count & 4 & 5 & conjectural & exploratory \\
large-field-operator-inequality & 5 & 6 & mixed & obstruction \\
lattice-gap-upper-bounds & 6 & 8 & mixed & theorem \\
lattice-truncation-uniform & 5 & 6 & mixed & obstruction \\
leibniz-continuity-records & 6 & 8 & conditional & reading \\
lieb-robinson-kogut-susskind & 5 & 7 & proved & theorem \\
local-detector-coincidences & 4 & 9 & proved & theorem \\
low-dimensional-mass-gap & 6 & 8 & mixed & theorem \\
magnetic-energy-identities & 4 & 8 & proved & obstruction \\
mark-cost-and-statistical-floor & 7 & 8 & proved & theorem \\
mass-gap-conditional-theorem & 4 & 4 & conditional & map \\
mass-gap-obligations-lattice & 4 & 7 & mixed & map \\
mass-gap-openings & 5 & 7 & mixed & map \\
mass-gap-position & 4 & 5 & mixed & map \\
mechanical-interference-action & 4 & 9 & proved & obstruction \\
millennium-problem-definitions & 2 & 8 & none & reading \\
minimax-composition & 4 & 8 & proved & calculation \\
moment-hierarchy-upper-bounds & 4 & 9 & conditional & theorem \\
necessity-unit-and-indeterminacy & 5 & 8 & proved & theorem \\
newton-NATP00385-audit & 3 & 7 & none & reading \\
newton-indeterminacy-routes & 6 & 8 & conditional & theorem \\
newton-insertion-action & 5 & 9 & proved & calculation \\
newton-mark-floor & 5 & 7 & mixed & theorem \\
newton-record-parallel-move & 3 & 8 & proved & calculation \\
ordered-beam-preparation & 1 & 7 & none & map \\
passive-threshold-events & 3 & 9 & proved & obstruction \\
physical-cut-speed & 3 & 9 & proved & obstruction \\
planck-gap-derivation & 6 & 6 & mixed & theorem \\
planck-gap-paper & 7 & 7 & mixed & theorem \\
planck-gap-probabilistic & 7 & 8 & proved & theorem \\
polyakov-average-gap-bound & 4 & 8 & proved & theorem \\
polygon-lift-phase & 6 & 9 & proved & theorem \\
position-preparation-ambiguity & 3 & 7 & proved & construction \\
principia-constant-force-action & 3 & 9 & proved & calculation \\
principia-fifth-postulate & 6 & 8 & proved & theorem \\
quantum-exclusion-premises & 3 & 8 & none & reading \\
radiation-noise-action-selection & 4 & 9 & proved & obstruction \\
reachable-cut-composition & 3 & 9 & proved & calculation \\
reasons-to-stop-as-research & 4 & 8 & mixed & map \\
receding-centre-area-audit & 3 & 9 & proved & calculation \\
reciprocal-coupling-normalization & 2 & 9 & proved & calculation \\
record-costs-disturbance & 6 & 8 & proved & theorem \\
record-costs-recoil & 6 & 8 & proved & theorem \\
record-distance-path-length & 6 & 8 & proved & theorem \\
refinement-composition-and-limit & 5 & 8 & proved & construction \\
refinement-results & 4 & 8 & mixed & map \\
relativistic-kepler-threshold & 4 & 9 & proved & calculation \\
reversible-generator-constraints & 3 & 9 & proved & reading \\
reversible-interaction-premise & 3 & 8 & conditional & reading \\
rivero-1998-conjecture-central-forces & 5 & 9 & proved & theorem \\
rotation-composition-universality & 4 & 8 & conditional & theorem \\
schur-error-ultraviolet & 5 & 4 & mixed & obstruction \\
score-constrained-ensemble & 5 & 8 & proved & construction \\
sed-closure-under-recording & 7 & 8 & mixed & obstruction \\
sed-zeta-radiation-link & 4 & 8 & conditional & reading \\
series-parallel-gauge-refinement & 7 & 8 & mixed & theorem \\
shared-readiness-chsh & 2 & 9 & proved & reading \\
shared-record-budget & 3 & 8 & proved & calculation \\
shared-resource-events & 2 & 9 & proved & calculation \\
single-calibration-fibres & 3 & 8 & proved & construction \\
small-field-step-decay-and-threshold & 5 & 6 & mixed & obstruction \\
small-field-step-gaussian & 4 & 7 & mixed & calculation \\
spin-action-patching & 3 & 9 & proved & calculation \\
stabilized-topology-action-scale & 3 & 9 & proved & calculation \\
static-composition-classics & 3 & 7 & none & map \\
stochastic-route-velocitas-ultima & 5 & 9 & conditional & theorem \\
strong-coupling-target-box & 5 & 5 & conditional & theorem \\
strong-coupling-threshold-explicit & 5 & 6 & mixed & calculation \\
strong-coupling-uniform-gap & 4 & 8 & proved & theorem \\
su2-midplane-order-t & 6 & 8 & conditional & calculation \\
su2-midplane-small-field & 6 & 7 & mixed & construction \\
su2-midpoint-exact & 5 & 9 & proved & calculation \\
su3-constants & 3 & 6 & mixed & map \\
sun-midpoint-centre & 6 & 8 & proved & theorem \\
superdeterminism-floor & 5 & 8 & conditional & theorem \\
susceptibility-gap & 3 & 9 & proved & theorem \\
tangent-groupoid-trajectories & 5 & 7 & mixed & reading \\
telegraph-return-bridge & 6 & 9 & proved & calculation \\
thermal-receiver-reliability & 4 & 9 & proved & calculation \\
thermodynamic-records-no-floor & 4 & 8 & proved & obstruction \\
three-body-cut-memory & 4 & 9 & proved & calculation \\
three-calibration-global-recovery & 4 & 7 & proved & theorem \\
three-continuum-limits & 4 & 8 & mixed & map \\
three-dimensional-gap-one-function & 4 & 7 & conditional & map \\
topological-sector-action-selection & 4 & 9 & proved & obstruction \\
torus-valley-potential & 5 & 8 & proved & calculation \\
two-calibration-branches & 4 & 8 & proved & construction \\
two-position-recovery & 3 & 9 & proved & calculation \\
two-regulator-audit & 3 & 8 & proved & map \\
typical-field-strength-window & 4 & 5 & superseded & obstruction \\
uv-halving-ir-confinement & 4 & 7 & none & map \\
villain-monopole-refinement & 6 & 9 & proved & theorem \\
weak-coupling-feshbach-reduction & 5 & 7 & conditional & theorem \\
what-would-unblock & 4 & 6 & none & map \\
wilson-strong-coupling-explicit & 5 & 8 & proved & calculation \\
zero-spacing-any-action & 6 & 8 & mixed & theorem \\
\end{longtable}

\subsection{5. Consequence for STATE}\label{consequence-for-state}

The two proof goals are unaffected: the Newton recording spine and the
refinement and gauge spine are the notes the scorer could follow to the
end. The pending work this audit creates is corrective: dated correction
boxes in the 24 notes of §3, the factor-2 fix in
\texttt{planck-gap-derivation}, the \(q=4\) column in
\texttt{strong-coupling-target-box}, the abstract of
\texttt{small-field-step-decay-and-threshold}, and one set of
strong-coupling thresholds with their labels in
\texttt{mass-gap-position} and \texttt{su3-constants}. STATE gains one
line pointing here.
\end{document}
